\documentclass[10pt,a4paper]{amsart}
\usepackage[margin=19mm]{geometry}
\usepackage{amsmath,amssymb,mathtools}
\usepackage[scr=rsfs,cal=euler]{mathalfa}
\usepackage{xfrac}
\usepackage[unicode,hidelinks,bookmarks=false]{hyperref}
\newcommand{\ie}{i.e.\ }
\newcommand{\cf}{cf.\ }
\newcommand{\resp}{respectively\ }
\newcommand{\Subsection}{\subsection}
\providecommand{\nobreakdash}{\nobreak\hbox{-}}
\setlength{\emergencystretch}{2em}
\multlinegap0pt
\usepackage{enumerate}
\usepackage{amssymb}
\usepackage{tikz}
\newtheorem{lemma}{Lemma}
\DeclareMathOperator{\Hom}{Hom}
\newcommand{\bbG}{\mathbb{G}}
\newcommand{\bbZ}{\mathbb{Z}}
\renewcommand{\le}{\leqslant}
\title{Exceptional Tannaka groups only arise\\ from cubic threefolds}
\author{Thomas Kr\"amer, Christian Lehn, Marco Maculan}
\date{Source: arXiv:2406.07401v3 (CC BY 4.0)}
\begin{document}
\maketitle

\noindent\emph{Source and licence.} Exceptional Tannaka groups only arise\\ from cubic threefolds. Version arXiv:2406.07401v3 (CC BY 4.0), distributed by arXiv under the Creative Commons Attribution 4.0 International licence, http://creativecommons.org/licenses/by/4.0/, as recorded in the arXiv licence metadata for this exact version and shown on the abstract page https://arxiv.org/abs/2406.07401v3. Full licence notice: \texttt{licenses/arxiv-2406.07401-CC-BY-4.0.txt}.\par\smallskip

\noindent\textit{Editorial scope.} Counted unit: the article's lemma lemma (source0.tex lines 1527--1529). The article's complete original proof of it is delivered with it (source0.tex lines 1531--1568, one \texttt{proof} environment of 38 lines). No mathematics is written, repaired or re-derived: the statement and the proof are the article's own lines, in the article's own notation, and no retained line is substituted or re-flowed.  No further source-local context is needed: every cross-reference of every retained line resolves inside the two delivered spans.  Nothing was omitted from any retained span, so the package carries no omission record.  No retained line contains a citation, so the wrapper carries no bibliography and no bibliography entry.  No retained line contains a figure environment, an image-inclusion command or drawing code, so no image is delivered and the package declares no asset dependency.  Editorial formatting only, in addition: an AMS \texttt{amsart} preamble that restates the article's own declarations and macros used by the retained lines, namely the environments \texttt{lemma} on separate counters, together with the standard packages those lines need.  The retained environments print in this document as Lemma 1 in order of appearance, whereas the article prints the same items with the numbering of its own full document.  The complete native source of the article as distributed in the arXiv e-print for this exact version is bundled unchanged as \texttt{sources/160647/source0.tex}.
\begin{lemma}
For any given~$(G, V)$, there exist up to~$G$-conjugacy only finitely many cocharacters~$\lambda$ that satisfy the property (H2). 
\end{lemma} 

\begin{proof} 
Fix a maximal torus~$T\subset G$. Any cocharacter is conjugate to one with values in this torus, so it suffices to consider cocharacters~$\lambda \in X_\ast(T) = \Hom(\bbG_m, T)$. The weight space decomposition for~$\lambda$ is then obtained as follows: Consider the weight space decomposition
%
\[
 V \;=\; \bigoplus_{\chi \in X^\ast(T)} V(\chi)
\]
%
for our maximal torus, where the action of the torus on the subspace~$V(\chi)\subset V$ is given by~$\chi \in X^\ast(T)=\Hom(T, \bbG)$. Then the weight~$n$ space for the cocharacter~$\lambda$ is the direct sum
%
\[
 V^n(\lambda) \;=\; \bigoplus_{\langle \lambda, \chi \rangle = n} V(\chi)
\]
%
where the sum is over all characters~$\chi \in X^\ast(T)$ with~$\langle \chi, \lambda \rangle = n$. To make use of this, let
%
\[
 S \;:=\; \{ \chi \in X^\ast(T) \mid V(\chi) \ne 0\}
\]
%
be the set of characters that enter in the given representation. Condition (H2) then implies
%
$|\langle \lambda, \chi \rangle| \le \ell < \dim V$
for all $\chi \in S$.
%
This condition leaves only finitely many choices of~$\lambda \in X_\ast(T)=\Hom(\bbG_m, T)$ in the cocharacter lattice, since in the dual character lattice the occurring weights~$\chi \in S$ span a subgroup
%
$
 \langle S \rangle_\bbZ \subset X^\ast(T)
$
%
of finite index: Indeed, if the subgroup spanned by those weights had smaller rank than~$X^\ast(T)$, then we could find~$\mu\in X_\ast(T)\smallsetminus \{0\}$ with~$\langle \mu, \chi\rangle = 0$ for all~$\chi \in S$. But then the image of
%
$
 \mu\colon \bbG_m \to G
$
%
would be a subgroup of positive dimension acting trivially on the representation~$V$ which is impossible since~$V$ has finite kernel.
\end{proof}

\end{document}
